Read in Gauge Theory Gravity: the metric $g$ below is the derived metric of the gravitational field $(\hb, \Omega)$ on flat spacetime, the manifold is the invertible region, and the tractor connection is the flat-space $\Cl(2,4)$ gauge field of Theorem \claimref{synthesis.gtg-tractor-gauge-field-reproduces-tractor-connection}; the index computation below is its reading in the splitting of $g$ (Fact \claimref{synthesis.imported-gtg-dictionary}).

\emph{(1)} In the parallel frame the transport of Fact \claimref{interior.imported-tractor-connection} reads, as in the proof of Proposition \claimref{synthesis.tractor-transport-weak-null-crossing}, $\sigma' = \mu(\gamma')$, $\mu' = -\gamma'^{\flat}\rho + \Sch(\gamma',\cdot)\,\sigma$, $\rho' = -\ip{\Sch(\gamma',\cdot)}{\mu}$. With $\Sch(\gamma',\cdot) = S\nu n$ and $n$ parallel, $w' = \ip{n}{\mu'} = -n(\gamma')\rho + S\nu\,Q(n)\,\sigma = -\nu\rho$, since $n$ is null, and $\rho' = -S\nu\ip{n}{\mu} = -\nu S w$. Differentiating, $w'' = -\nu\rho' = \nu^2 S w$, and $\rho = -w'/\nu$. The check verifies the two row-vector identities $wA = -\nu\rho$, $\rho A = -\nu S w$ on five frames (timelike and null $\gamma'$) and five values of $S$.

\emph{(2)} Work on one side, with $t = |s - s_c|$. (a) If $S$ is integrable, $A$ is integrable and the whole transport has a limit (Proposition \claimref{synthesis.tractor-transport-weak-null-crossing}(3)). (b) If $\int t|S| < \infty$, write $w(s) = w(s_0) + w'(s_0)(s - s_0) + \int_{s_0}^{s}(s - \tau)\nu^2 S(\tau) w(\tau)\,d\tau$; the kernel is bounded by $\nu^2|s_c - \tau|\,|S(\tau)|$, integrable, so Picard iteration gives a bounded solution with a limit $w(s_c)$, for every datum. Then $w'(s) = w'(s_0) + \nu^2\int_{s_0}^{s} S w$; if $w(s_c) \neq 0$, $w$ keeps the sign of $w(s_c)$ near $s_c$ with $|w| \ge |w(s_c)|/2$, and the fixed sign of $S$ with $\int|S| = \infty$ make $|w'|$, hence $|\rho|$, unbounded. The map from data at $s_0$ to $w(s_c)$ is linear and nonzero (the positive solution of the stored theorem has $w(s_c) > 0$), so $w(s_c) = 0$ is a line. (c) Suppose $\int t|S| = \infty$ on the side. If some solution had a limit $L \neq 0$, say $L > 0$, then $w \ge L/2$ near $0$ and, in the variable $t$ (where $d^2w/dt^2 = \nu^2 S w$ as well), $w(t) = w(t_0) - \tfrac{dw}{dt}(t_0)(t_0 - t) + \int_t^{t_0}(\tau - t)\nu^2 S(\tau) w(\tau)\,d\tau$, whose last term has the fixed sign of $S$ and modulus at least $\nu^2(L/2)\int_t^{t_0}(\tau - t)|S|$, unbounded as $t \to 0$ by monotone convergence: a contradiction. So no $w$, and no gauge solution $y$, has a finite nonzero limit. For $S \ge 0$ take the solution with $w(t_0) = 1$, $dw/dt(t_0) = 0$: as long as $w \ge 1$ on $[t, t_0]$, $w(t) \ge 1 + \nu^2\int_t^{t_0}(\tau - t)S(\tau)\,d\tau$, which keeps $w \ge 1$ (continuity argument) and tends to infinity. So an unbounded solution exists; the bounded solutions form a proper subspace of the two-dimensional solution space, at most a line. A change of splitting with bounded parameters maps bounded tractors to bounded ones, so the unbounded $\ip{n}{\mu}$ is not an artefact of the splitting of $g$ within that class. The sign statement: with $\mu_a = \nabla_a\sigma$ for a parallel tractor, the transport gives $\sigma'' = \Sch(\gamma',\gamma')\sigma$ along a null geodesic, against $\sigma'' = -\mathrm{P}^{\mathrm{BEG}}(\gamma',\gamma')\sigma$ in the Bailey-Eastwood-Gover convention, consistently with Lemma \claimref{interior.sign-fixed-internally}'s naming $\Ric = \Lambda g$ for de Sitter in signature $(+,-,-,-)$; so focusing matter, $\mathrm{P}^{\mathrm{BEG}}(l, l) \ge 0$, has $S \le 0$.

\emph{(3)} For $S = \pm Ct^{-2}$: $t^{e}$ solves $d^2y/dt^2 = \pm\nu^2Ct^{-2}y$ iff $e(e - 1) = \pm\nu^2 C$, and $\int_0^{t_0} t|S|\,dt = C\int_0^{t_0}dt/t = \infty$ (the measure $t|S|\,dt = C\,dt/t$ is invariant under $t \mapsto t/2$, so each dyadic shell carries the same positive amount). For $+$ the roots are $k > 1$ and $1 - k < 0$. For $-$ the roots satisfy $e_1 + e_2 = 1$, $e_1e_2 = \nu^2 C > 0$: both real in $(0, 1)$ if $\nu^2 C \le \tfrac14$ (a double root $\tfrac12$ with solutions $t^{1/2}$, $t^{1/2}\log t$ at equality), and $\tfrac12 \pm I_{4}\omega$ (numbers read in $\mathrm{span}\{1, I_{4}\}$), $\omega^2 = \nu^2 C - \tfrac14$, otherwise, with real solutions $t^{1/2}\cos(\omega\log t)$, $t^{1/2}\sin(\omega\log t)$. In every focusing case all solutions tend to $0$, and their derivatives grow at most like $t^{\mathrm{Re}\,e - 1}$, integrable; in the complex case the zeros accumulate at $0$. For a real root $e$ and $y = t^e$, with $d/ds = -d/dt$ on the side $s < s_c$, $f = y'/(\nu y) = -e/(\nu t)$ and $f' = -e/(\nu t^2) = \nu(S - f^2)$ because $e^2 - e = \nu^2(\pm C)$; so the stored theorem's algebraic identity $(1 - fN_n)A(1 + fN_n) - f'N_n = E_{\gamma'} + fD$ holds (it uses only $N_n^2 = 0$ and $N_nE_{\gamma'}N_n = \nu N_n$; the check verifies it at rational values and that it fails when $f'$ violates the Riccati equation). $D = [E_{\gamma'}, N_n]$ acts by $\nu$ on $\sigma$, $-\nu$ on $\rho$ and as a boost generator on $\mu$; the check verifies $D^3 = \nu^2 D$ and that $D$, $D - \nu$, $D + \nu$ all have rank four on five frames, so $D$ is diagonalisable with eigenvalues $0, \nu, -\nu$, each twice. Hence $d\hat x/dt = \bigl(\tfrac{e}{\nu t}D - E_{\gamma'}\bigr)\hat x$ has a regular singular point with residue $\tfrac{e}{\nu}D \in \mathfrak g_0$ of eigenvalues $0, \pm e$.

\emph{(4)} For $|S| = Ct^{2\beta - 2}$ (either sign; only $|S|$ enters (2)(a), (b)): $\int_0^1 t|S|\,dt = C/(2\beta)$ for every $\beta > 0$, and $\int_\epsilon^1 S = C(1 - \epsilon^{2\beta - 1})/(2\beta - 1)$ is bounded as $\epsilon \to 0$ iff $\beta > \tfrac12$ (at $\beta = \tfrac12$ it is $C\log(1/\epsilon)$). The Frobenius series $y = \sum_j a_j t^{2\beta j}$, $a_0 = 1$, $a_j = \nu^2 C a_{j-1}/\bigl(2\beta j(2\beta j - 1)\bigr)$ (when $2\beta j \neq 1$ for all $j$) solves $d^2y/dt^2 = \nu^2 C t^{2\beta - 2} y$ (for $S \le 0$ replace $C$ by $-C$), converging for all $t$ by the factorial growth of the denominators; $y(0) = 1$ and $dy/dt = \nu^2 C t^{2\beta - 1}/(2\beta - 1) + O(t^{4\beta - 1})$, giving the stated rate of $|f|$. The check verifies the series to six terms exactly (only the truncation term survives), its leading derivative, and the partial integrals at $\epsilon = 8^{-m}$ for $\beta = \tfrac13$ (growing like $2^m$) and $\beta = \tfrac23$ (bounded by $3$). For $|S| = Ct^{-2}(\log 1/t)^{-q}$ on $t < 1/\mathrm{e}$, the substitution $x = \log(1/t)$ gives $t|S|\,dt = -Cx^{-q}\,dx$, so the Tipler integral is $C\int_1^\infty x^{-q}\,dx$, finite iff $q > 1$, while $|S|$, $t^{-2}$ times a slowly varying factor, is never integrable. Combining with (2) gives the stated cases.
